\documentclass[a4paper,12pt]{article}

% Pacchetti comuni
\usepackage[utf8]{inputenc}    % Codifica UTF-8
\usepackage{amsmath}            % Simboli matematici
\usepackage{amsfonts}           % Font matematici
\usepackage{amssymb}            % Simboli matematici
\usepackage{graphicx}           % Gestione delle immagini
\usepackage{hyperref}           % Collegamenti ipertestuali

\title{Appunti della lez}
\author{Nicolo' Giuliani}
\date{\today}

\begin{document}
\section{Ripasso}
L'insieme delle matrici $m * n$ e' un anello.
\\Abbiamo 2 operazioni:
\\1) Somma e' commutativa, associativa, esiste elemento neutro ed esisto l'opposto.
\\2) Prodotto righe per colonne non e' commutativo se $m \geq n$. E' associativo, vale la prorpieta' distributiva.
\\Esiste l'elemento neutro:
\begin{align*}
    \begin{pmatrix}
        1 & 0 &0 & 0 \\
        0 & 1 & 0 & 0 \\
        0 & 0 & 1 & 0 \\
        0 & 0 & 0 & 1
    \end{pmatrix}
    =I \text{ (matrice identita')}
\end{align*}
\\$AI=IA \, \forall A \in M_{m,n}(\mathbb{R})$
\\$m$ si dice l'ordine della matrice.

\section{Determinante}
\textit{Il determinante se non si tratta di matrici quadrate non esiste.}
Il determinante e' una funzione $M_{n,n}(\mathbb{R})\to \mathbb{R}$ con varie prorpieta' \textit{(vedere libro)}.
\begin{align*}
    A= \begin{pmatrix}
        1 & 0 & 5 \\
        3 & -1 & 4 \\
        0 & 2 & 7
    \end{pmatrix}
    \\ R_2=\begin{pmatrix}
        3 & -1 & 4
    \end{pmatrix}
    =(1,0,-7)+(2,-1,11) 
    \\\text{ (ovvero la somma di due matrici)}
    \\\text{det A}=det\begin{pmatrix}
        1 & 0 & 5 \\ 1 & 0 & -7 \\ 0 & 2 & 7
    \end{pmatrix} + 
    det\begin{pmatrix}
        1 & 0 & 5 \\ 2 & -1 & 11 \\ 0 & 2 & 7
    \end{pmatrix}
\end{align*}
\begin{align*}
    A=\begin{pmatrix}
        3 & -1 & 8 \\ 6 & 1 & 5 \\ 9 & -3 & 12
    \end{pmatrix}
    \\ \begin{pmatrix}
        9 & -3 & 12
    \end{pmatrix}
    =3\begin{pmatrix}
        3 & -1 & 4
    \end{pmatrix}
    \\\text{det A}=3det\begin{pmatrix}
        \\
        \\
        1 & 2 & 7
    \end{pmatrix}
\end{align*}
\begin{align*}
    \text{b) }
    \begin{pmatrix}
        R_1 \\ R_2 \\ \vdots \\ R_n
    \end{pmatrix} = A
    \\ \begin{pmatrix}
        R_1 \\
        R_2 \\
        R_i+3R_1+7R_2-9R_n \\
        R_n
    \end{pmatrix} = B
    \\ \text{det A = det B}
\end{align*}
\begin{align*}
    \text{c) una matrice triangola superiore e' del tipo:}
    \\ \begin{pmatrix}
        x & x & x & x \\
        0 & x & x & x \\
        0 & 0 & 0 & x \\
        0 & 0 & 0 & 0
    \end{pmatrix}
    \\\text{Una matrice triangolare inferiore e':}
    \\\begin{pmatrix}
        x & 0 & 0 \\
        x & x & 0 \\
        x & x & x
    \end{pmatrix}
\end{align*}
Esempi:
\begin{align}
    det\begin{pmatrix}
        3 & 9 & 7 \\
        0 & 1 & 4 \\
        0 & 0 -1
    \end{pmatrix}=3*1*(-1)=-3
    \\ det\begin{pmatrix}
        3 & 2 & 0 \\
        0 & 0 & 1 \\
        0 & 0 & 0
    \end{pmatrix}=0
\end{align}
Una matrice a scala e' sicuramente triangolare. \\
\pagebreak
Grazie alle proprieta' precedenti possiamo calcolare il determinante con l'algoritmo di Gauss.
\begin{align*}
    A=\begin{pmatrix}
        1 & 2 & 0 \\-1&3&1\\0&1&-2
    \end{pmatrix}
    \\R_2+R_1=\begin{pmatrix}
        1&2&0\\0&5&1\\0&1&-2
    \end{pmatrix}=A_1 \\
    \text{det A = det }A_1
    \\ R_3 \iff R_2 \to \begin{pmatrix}
        1&2&0\\0&1&-2\\0&5&1
    \end{pmatrix}=A_2
    \\\begin{pmatrix}
        1&2&0\\0&1&-2\\0&0&-11
    \end{pmatrix}=A_3 \\
    \text{det }A_2 \text{=det }A_3
    \\\text{det }A_3=1*1*11=11
    \\\text{det A}=\text{det }A_1=-\text{det }A_2=-\text{det }A_3=-11 \\
\end{align*}
Metodi per il cacolo del determinante
\begin{align*}
    2*2: \\
    A=\begin{pmatrix}
        a&b\\c&d
    \end{pmatrix} \\
    \text{det }A=ad-bc \\
    3*3: \\
    \begin{pmatrix}
        a_{11}&a_{12}&a_{13}\\a_{21}&a_{22}&a_{23}\\a_{31}&a_{32}&A_{33}
    \end{pmatrix} \\
    \text{det} = a_{11}a_{22}a_{33} + a_{12}a_{23}a_{31} + a_{13}a_{21}a_{32} - a_{13}a_{22}a_{31} - a_{12}a_{21}a_{33} - a_{11}a_{23}a_{32}
\end{align*}
Esempio:
\begin{align*}
    A=\begin{pmatrix}
        1&2&0&1&2\\-1&3&1&-1&3\\0&1&-2&0&1
    \end{pmatrix} \\
    \text{consideriamo }\begin{pmatrix}
        1&2&0\\-1&3&1\\0&1&-2
    \end{pmatrix} \\
    \text{det A}=1*3*(-2)+2*1*0+0*(-1)*1=-11
\end{align*}
Metodo di Laplace:
\\\textit{sviluppo secondo la prima riga}
\begin{align*}
    T_{i,j}=(-1)^{i+j}det(A_{ij})
    \\ A=\begin{pmatrix}
        a_{11}&a_{12}&\dots&a_{1m} \\
        a_{21}&a_{22}&\dots&a_{2m}\\
        \vdots \\
        a_{m1}&a_{m2}&\dots&a_{mm}
    \end{pmatrix}
    \\detA=a_{i1}T_{i1}+a_{i2}T_{i2}+\dots+a_{im}T_{im}
\end{align*}
Esempio:
\begin{align*}
    A=\begin{pmatrix}
        a&b\\c&d
    \end{pmatrix}
    \\detA=aT_{11}+bT_{12}
    \\=a(-1)^{1+1}det(d)+b(-1)^{1+2}det(c)=ad-bc
\end{align*}
Una matrice $A\in M_{n,n}(\mathbb{R})$ si dice invertibile se esiste una matrice B tale che $AB=BA=I$
\\B si dice la matrice inversa di A e si indica con $A^{-1}$.
\subsection{Proposizione}
Una matrice $A\in M_{n,n}(\mathbb{R})$ e' invertibile solo se $detA\neq0$.
\subsection{Teorema di Binet}
Se $A,B\in M_{n,n}(\mathbb{R})$ allora:
\\det AB = det A * det B \\
\textit{RICORDA:} $det (A+B) \neq det A + det B$
\subsection{Proposizione}
Se $A\in M_{n,n}(\mathbb{R})$ ha una riga o una colonna uguale ad 0 allora det(A)=0.
\subsection{Proposizione}
A e' invertibile se $det A \neq 0$.
\\Dim:
\begin{align*}
    B=A^{-1} \\
    AB=I \to det(I)=det(AB)=det(A)det(B)
    \\\text{Ricorda: det(I)=1}
    \\\text{quindi } det(A)det(B)=1 \text{ allora } detA\neq 0
\end{align*}
\begin{align*}
    \text{Sia } detA\neq0
    B=b_{ij}=\frac{1}{detA}(-1)^{i+j}det(A_{ij})
    \\ =\frac{1}{detA}T_{ij}
\end{align*}
Si dimostra che $B=A^{-1}$

\end{document}
